\chapter{The auction: who gets what}\label{ch:auction}

\begin{chapterplan}
The question of this chapter: when the incumbent is stronger, how does the sale divide the challenger's value
between the challenger and the target?

After this chapter, you can do these things:
\begin{itemize}
\item Explain why each prepared bidder bids its value.
\item For any incumbent value and any challenger value, write what the target receives and what the challenger
keeps, and show that the two add up to the challenger's value.
\item Derive the gross acquisition profit \(g_j\) and the target-payoff spread \(\DT\) as integrals over levels
of value.
\item State and prove Proposition~1, with its conditions for a strict change, and explain its part~(c).
\item Explain the class-gap identity: a stronger incumbent moves the class value gap from the challenger's profit
to the target's proceeds.
\end{itemize}
\end{chapterplan}

% ====================================================================
\section{The story in words}\label{sec:auction-story}

The sale is a cash second-price auction with a reserve \(p>0\). The highest admissible bid wins. The winner pays
the larger of the reserve and the highest competing admissible bid. Section~\ref{sec:auction-bids} shows that
each prepared bidder bids its own value. So, when at least one value meets the reserve, the price is the lower of
the two values, or the reserve when the reserve is higher.

Take a challenger that prepared and so learned its value \(\theta\), with \(\theta\ge p\). The incumbent has the
value \(R\). Two outcomes matter:
\begin{itemize}
\item If \(R<\theta\), the challenger wins and pays \(\max\{p,R\}\). The target's proceeds then follow the
incumbent's value.
\item If \(R>\theta\), the incumbent wins and pays \(\theta\), the challenger's bid. The target's proceeds then
follow the challenger's value.
\end{itemize}

First look at the challenger. Suppose that the incumbent's value \(R\) is a little higher. Then no more challenger
values beat it, and a challenger that still wins pays at least as much. Neither change raises what the challenger
keeps. A stronger incumbent makes high values of \(R\) more likely, so it weakly lowers the challenger's expected
profit. The paper calls this the direct deterrence force of a strong rival: it lowers the return to preparation.

Next look at the target. A weak incumbent usually loses, and the target then receives \(\max\{p,R\}\). This
amount does not depend on the challenger's value at all. A strong incumbent often wins and pays \(\theta\), and
the target's proceeds then equal the challenger's value. So a stronger incumbent makes the target's proceeds
depend more on how good the challenger is.

The paper keeps the deterrence force and adds a channel that starts from this second effect. The challenger comes
from a high class or a low class, and the
high class tends to have higher values. The expected difference between the values of the two classes is the
\emph{class value gap}. The auction gives each part of this gap either to the target or to the challenger. The
incumbent's value decides the division, and a stronger incumbent moves parts of the gap to the target. The paper
calls these parts \emph{slices}.

\begin{keyidea}
The incumbent's value works as a boundary line on the challenger's value. Above the reserve, the part of the
challenger's value below the line goes to the target, and the part above the line stays with the challenger. A
stronger incumbent raises the line.
\end{keyidea}

\begin{analogy}[a price clock]
Picture an auctioneer who raises a price step by step, from the reserve upward. Each bidder stays in while the
price is below its value, and it leaves when the price reaches its value. The sale stops when one bidder leaves,
and the other bidder buys at that price. When each bidder knows its own value, this clock gives the same winner
and the same price as the sealed second-price auction. Each step of the clock, from a level \(y\) to \(y+dy\), is
a slice. The target collects the slice when both bidders are still in at the level \(y\).

Where the analogy breaks: in the model, the bids are sealed and are made at the same time, so no bidder sees the
other leave. The clock is only a picture of the outcome.
\end{analogy}

Two extreme incumbents make the picture sharp. For this paragraph, assume that every challenger value meets the
reserve.
\begin{itemize}
\item An incumbent that never meets the reserve never bids. The challenger always wins and pays \(p\). The target
receives \(p\) from either class, so the class does not matter to the target. The whole class value gap stays
with the challenger.
\item An incumbent whose value is always above every challenger value always wins and pays \(\theta\). The
challenger keeps nothing, and the target collects the whole challenger value. The whole class value gap goes to
the target.
\end{itemize}
A strengthening of the incumbent moves the sale from the first extreme toward the second.
Section~\ref{sec:auction-gap} turns this picture into an exact identity.

% ====================================================================
\section{The set-up}\label{sec:auction-setup}

This section recalls the parts of the model that the auction uses. Chapter~\ref{ch:model} gives the full model.

\subsection*{The sale rule}

Before trade, the seller commits to a cash second-price auction with a reserve \(p>0\). The highest admissible
bidder acquires the target. It pays the larger of the reserve and the highest competing admissible bid. A bid
equal to the reserve is admissible, and there is no sale without an admissible bid. A value \emph{meets the
reserve} when it is at least \(p\). The paper then calls the value, and the bidder, admissible.

\subsection*{The two bidders}

\begin{itemize}
\item The incumbent prepared before the sale. It knows its value \(R\) before it bids. The value \(R\) has a
continuous distribution function \(F\) on \([0,\bar r]\), and \(F(y)=1\) for \(y\ge\bar r\).
\item The challenger belongs to the class \(j\in\{H,L\}\). Given the class, its value \(\theta\) has a
distribution function \(G_j\) on \([0,\infty)\) with a finite mean. The challenger learns \(\theta\) only when it
prepares.
\item The pair \((j,\theta)\) and the incumbent's value \(R\) are independent.
\end{itemize}

Continuity of \(F\) means that \(\Pr(R=y)=0\) for every \(y\). So the strict and the weak inequality give the
same probability (Section~\ref{sec:toolkit-cdf}):
\[
\Pr(R<y)=\Pr(R\le y)=F(y),\qquad \Pr(R>y)=\Pr(R\ge y)=1-F(y).
\]
In particular, a tie \(R=\theta\) has probability zero.

The class distributions \(G_j\) can have atoms. An atom at \(y\) means that the value equals \(y\) with a
positive probability, so \(G_j\) jumps at \(y\). Write \(G_j(y^-)\) for the left limit: the limit of \(G_j(x)\) as
\(x\) rises to \(y\). Then the four probabilities about one level are (Section~\ref{sec:toolkit-cdf}):
\begin{align*}
\Pr\nolimits_j(\theta\le y)&=G_j(y), & \Pr\nolimits_j(\theta<y)&=G_j(y^-),\\
\Pr\nolimits_j(\theta>y)&=1-G_j(y), & \Pr\nolimits_j(\theta\ge y)&=1-G_j(y^-).
\end{align*}
The two forms differ only at an atom. The probability that a class-\(j\) value falls below the reserve is
\(G_j(p^-)\), and the probability that it meets the reserve is \(1-G_j(p^-)\).

The paper compares the classes level by level with
\[
D(y)=G_L(y)-G_H(y)=\bigl[1-G_H(y)\bigr]-\bigl[1-G_L(y)\bigr]=\Pr\nolimits_H(\theta>y)-\Pr\nolimits_L(\theta>y).
\]
So \(D(y)\) says how much more likely a high-class value is than a low-class value to lie above the level \(y\).

\subsection*{The class ranking and equal admissibility}

The paper's main condition on the classes is the \emph{class ranking}:
\begin{equation}\label{eq:auction-ranking}
G_H(y)\le G_L(y)\quad\text{for every }y,\qquad G_H(p^-)=G_L(p^-)=0.
\end{equation}
The first part is first-order stochastic dominance (Section~\ref{sec:toolkit-fosd}): at every level, a high-class
value is at least as likely as a low-class value to lie above it. The second part says that both class supports
lie above the reserve, so every value meets it. In terms of \(D\), the class ranking implies \(D(y)\ge0\) for
every \(y\) and \(D(p^-)=0\). These two facts are weaker than the class ranking: they do not need both classes to
meet the reserve with probability one.

A weaker condition appears in Section~\ref{sec:auction-weakest}. The classes have \emph{equal admissibility} when
\(G_H(p^-)=G_L(p^-)\): the two classes fall below the reserve with the same probability. The class ranking
implies equal admissibility, with both probabilities equal to zero.

\subsection*{The payoff objects}

\begin{itemize}
\item \(t_0\) is the expected target proceeds without entry.
\item For a known challenger value \(\theta\), \(t(\theta)\) is the expected target proceeds after the
challenger prepares, and \(g(\theta)\) is the challenger's expected gross acquisition profit. Both expectations
are over the incumbent's value \(R\).
\item The class objects average over the class distribution: \(t_j=\E_j[t(\theta)]\) and
\(g_j=\E_j[g(\theta)]\) for \(j\in\{H,L\}\).
\item The target-payoff spread is \(\DT=t_H-t_L\).
\item At a belief \(\mu=\Pr(j=H)\), the challenger's expected gross profit is
\(B(\mu)=g_L+\mu\,(g_H-g_L)=(1-\mu)\,g_L+\mu\,g_H\).
\end{itemize}
When the incumbent's distribution matters, write \(g_j(F)\), \(\DT(F)\) and \(B_F\). The positive part of a
number \(x\) is \((x)_+=\max\{x,0\}\).

Without a challenger, the incumbent is the only possible bidder. It bids when \(R\ge p\), and then it pays
\(\max\{p,0\}=p\), since there is no competing bid. So:
\begin{align*}
t_0&=p\cdot\Pr(R\ge p)+0\cdot\Pr(R<p) \why{the incumbent buys at \(p\), or no sale}\\
&=p\,\bigl[1-\Pr(R<p)\bigr] \why{complement rule}\\
&=p\,\bigl[1-F(p)\bigr] \why{\(F\) is continuous}
\end{align*}

A \emph{strengthening} of the incumbent replaces \(F\) by a continuous \(\tilde F\) with \(\tilde F(y)\le F(y)\)
for every \(y\). At every level \(y\), the stronger incumbent is at or below \(y\) with at most the same
probability. This is a first-order stochastic shift (Section~\ref{sec:toolkit-fosd}).

\begin{watchout}
The probability that a class-\(j\) value fails the reserve is \(G_j(p^-)\), with the left limit. A bid equal to
the reserve is admissible, so a value exactly at \(p\) meets the reserve. When class \(j\) has an atom at \(p\),
the number \(G_j(p)\) counts that atom and so overstates the share of values that fail the reserve.
\end{watchout}

Bids turn out to be truthful (Section~\ref{sec:auction-bids}). So every object above depends on the incumbent's
distribution, the classes and the reserve, and on no belief. Prices and trade play no role in this chapter.

% ====================================================================
\section{Bids: why each bidder bids its value}\label{sec:auction-bids}

The paper's equilibrium definition lists truthful bids as one condition, (E6). The appendix then shows that a
bid equal to the bidder's value is weakly dominant, so the condition holds in every equilibrium tuple. A bid
is \emph{weakly dominant} when, against every possible competing bid, it gives at least as much as any other bid
(Section~\ref{sec:toolkit-games}). A bidder needs no belief about its rival to choose such a bid.

Fix one bidder with the value \(\theta\). Let \(Y\ge0\) be the highest competing admissible bid, with \(Y=0\) when
there is none. The rule gives the bidder a simple choice:
\begin{itemize}
\item If it wins with an admissible bid, it pays \(\max\{p,Y\}\). This price does not depend on its own bid.
\item If it loses, or makes no bid, it pays nothing and gets nothing.
\end{itemize}
So every bid yields one of two payoffs: \(\theta-\max\{p,Y\}\) or zero. The bid only decides which of the two
payoffs the bidder gets.

\subsection*{A value at or above the reserve}

Suppose \(\theta\ge p\), and consider the bid \(\theta\). It is admissible. There are three cases.
\begin{enumerate}
\item \(\theta>\max\{p,Y\}\). Then \(\theta>Y\), so the bid \(\theta\) wins. It earns \(\theta-\max\{p,Y\}>0\).
\item \(\theta<\max\{p,Y\}\). Since \(\theta\ge p\), the larger of \(p\) and \(Y\) must be \(Y\), so
\(Y>\theta\). The bid \(\theta\) loses and earns zero.
\item \(\theta=\max\{p,Y\}\). If the bid wins, it pays \(\theta\) and earns zero. If it loses, it also earns
zero.
\end{enumerate}
In every case, the bid \(\theta\) earns
\[
\max\bigl\{\theta-\max\{p,Y\},\,0\bigr\}=\bigl(\theta-\max\{p,Y\}\bigr)_+,
\]
the larger of the two possible payoffs. No bid can earn more than the larger of the two. So the bid \(\theta\) is
weakly dominant. A bid below the reserve is not admissible, so it has the same effect as no bid and earns zero.
Check yourself~\ref{q:auction-overbid} compares the bid \(\theta\) with a bid above the value, range by range.

\subsection*{A value below the reserve}

Suppose \(\theta<p\). Any bid that wins pays \(\max\{p,Y\}\ge p>\theta\). So a win earns
\(\theta-\max\{p,Y\}<0\). No bid earns more than zero, and the choice of no bid earns zero against every \(Y\).
So the choice of no bid is weakly dominant. This is the paper's rule: a prepared challenger whose value is below the reserve does
not bid.

\subsection*{What the bids give the rest of the chapter}

\begin{itemize}
\item Truthful bids, condition (E6) of the paper's equilibrium definition, hold in every equilibrium tuple. No
payoff in this chapter rests on a belief.
\item The challenger faces the incumbent. The incumbent bids \(R\) when \(R\ge p\) and makes no bid when
\(R<p\). So \(Y=R\) when \(R\ge p\), and \(Y=0\) when \(R<p\). In both cases, a challenger that wins pays
\(\max\{p,R\}\):
\begin{align*}
R\ge p:\quad \max\{p,Y\}&=\max\{p,R\} \why{\(Y=R\)}\\
R<p:\quad \max\{p,Y\}&=\max\{p,0\}=p \why{\(Y=0\)}\\
&=\max\{p,R\} \why{\(R<p\)}
\end{align*}
\end{itemize}

\begin{intuition}[why the value is the right bid]
The rule separates the price from the decision to buy. The other bidder sets the price, and your own bid only
decides whether you buy at that price. You want to buy exactly when your value is above the price. A bid equal to your value buys in
exactly those cases.
\end{intuition}

\begin{checkyourself}
\label{q:auction-overbid}
A bidder with the value \(\theta\ge p\) considers the bid \(b>\theta\). Go through the competing bids \(Y<\theta\),
\(\theta\le Y\le b\) and \(Y>b\). In each range, compare the payoff of the bid \(b\) with the payoff of the bid
\(\theta\). Where is the bid \(b\) strictly worse?
\end{checkyourself}

% ====================================================================
\section{What each party receives at fixed values}\label{sec:auction-outcomes}

Fix the incumbent's value \(R\) and the value \(\theta\) of a prepared challenger. This section finds who wins,
what the target receives and what the challenger keeps. The expectations come in the later sections.

\subsection*{A challenger value at or above the reserve}

Start with \(\theta\ge p\). Table~\ref{tab:auction-cases} lists the three cases.

\begin{table}[htbp]
\centering\small
\caption{A prepared challenger with the value \(\theta\ge p\) against an incumbent with the value \(R\).}
\label{tab:auction-cases}
\begin{tabular}{@{}lllll@{}}
\toprule
Case & Incumbent's bid & Winner & Target receives & Challenger keeps\\
\midrule
\(R<p\le\theta\) & none & challenger & \(p\) & \(\theta-p\)\\
\(p\le R<\theta\) & \(R\) & challenger & \(R\) & \(\theta-R\)\\
\(p\le\theta<R\) & \(R\) & incumbent & \(\theta\) & \(0\)\\
\bottomrule
\end{tabular}
\end{table}

\begin{itemize}
\item \(R<p\le\theta\). The incumbent's value is below the reserve, so it does not bid. The challenger is the
only admissible bidder. It wins and pays \(\max\{p,0\}=p\).
\item \(p\le R<\theta\). Both bidders bid their values. The challenger has the higher bid. It pays the competing
bid \(R\), which is at least \(p\).
\item \(p\le\theta<R\). The incumbent has the higher bid. It pays the challenger's bid \(\theta\), which is at
least \(p\).
\end{itemize}

Two formulas cover all three rows. For \(\theta\ge p\):
\[
\text{the target receives }\max\{p,\min(R,\theta)\},\qquad
\text{the challenger keeps }\bigl(\theta-\max\{p,R\}\bigr)_+.
\]

\begin{derivation}[the two formulas reproduce the table]
\begin{align*}
R<p\le\theta:\quad \max\{p,\min(R,\theta)\}&=\max\{p,R\}=p \why{\(\min(R,\theta)=R\); \(R<p\)}\\
\bigl(\theta-\max\{p,R\}\bigr)_+&=(\theta-p)_+=\theta-p \why{\(\max\{p,R\}=p\); \(\theta\ge p\)}\\
p\le R<\theta:\quad \max\{p,\min(R,\theta)\}&=\max\{p,R\}=R \why{\(R<\theta\); \(R\ge p\)}\\
\bigl(\theta-\max\{p,R\}\bigr)_+&=(\theta-R)_+=\theta-R \why{\(R\ge p\); \(R<\theta\)}\\
p\le\theta<R:\quad \max\{p,\min(R,\theta)\}&=\max\{p,\theta\}=\theta \why{\(\theta<R\); \(\theta\ge p\)}\\
\bigl(\theta-\max\{p,R\}\bigr)_+&=(\theta-R)_+=0 \why{\(R>\theta\ge p\); \(\theta-R<0\)}
\end{align*}
\end{derivation}

\emph{Ties.} At \(R=\theta\), whichever bidder wins pays \(\theta\). So the target receives \(\theta\), and the
winner earns zero. The two formulas give \(\max\{p,\theta\}=\theta\) and \((\theta-\theta)_+=0\). So the
formulas hold under any tie rule. Since \(F\) is continuous, the tie has probability zero anyway.

\subsection*{A challenger value below the reserve}

Suppose \(\theta<p\). The challenger does not bid (Section~\ref{sec:auction-bids}). If \(R\ge p\), the incumbent
is the only admissible bidder and buys at \(p\). If \(R<p\), there is no sale. So the target receives
\(p\,\1\{R\ge p\}\), and the challenger keeps nothing. This is the same outcome as without entry.

\begin{figure}[htbp]
\centering
\begin{tikzpicture}[x=1cm,y=1cm,every node/.style={font=\small}]
  % row (i): R < p <= theta
  \node[anchor=east] at (-0.4,0) {\(R<p\le\theta\)};
  \draw[rulegrey] (0,0) -- (10,0);
  \draw[accent,line width=3pt] (0,0) -- (4,0);
  \draw[notegrey!45,line width=3pt] (4,0) -- (8.5,0);
  \foreach \x/\t in {0/0,4/p,8.5/\theta} {\draw (\x,0.18) -- (\x,-0.18) node[below] {\(\t\)};}
  \draw[notegrey] (2,0.18) -- (2,-0.18) node[below] {\(R\)};
  % row (ii): p <= R < theta
  \node[anchor=east] at (-0.4,-1.25) {\(p\le R<\theta\)};
  \draw[rulegrey] (0,-1.25) -- (10,-1.25);
  \draw[accent,line width=3pt] (0,-1.25) -- (5.5,-1.25);
  \draw[notegrey!45,line width=3pt] (5.5,-1.25) -- (8.5,-1.25);
  \foreach \x/\t in {0/0,2.5/p,5.5/R,8.5/\theta} {\draw (\x,-1.07) -- (\x,-1.43) node[below] {\(\t\)};}
  % row (iii): p <= theta < R
  \node[anchor=east] at (-0.4,-2.5) {\(p\le\theta<R\)};
  \draw[rulegrey] (0,-2.5) -- (10,-2.5);
  \draw[accent,line width=3pt] (0,-2.5) -- (5.5,-2.5);
  \foreach \x/\t in {0/0,2.5/p,5.5/\theta,8.5/R} {\draw (\x,-2.32) -- (\x,-2.68) node[below] {\(\t\)};}
  % legend
  \draw[accent,line width=3pt] (0,-3.6) -- (0.8,-3.6);
  \node[anchor=west] at (0.9,-3.6) {target receives};
  \draw[notegrey!45,line width=3pt] (4.2,-3.6) -- (5.0,-3.6);
  \node[anchor=west] at (5.1,-3.6) {challenger keeps};
\end{tikzpicture}
\caption{The split of an admissible challenger value at a fixed incumbent value. In each row, the dark part of the
line, from zero to \(\max\{p,\min(R,\theta)\}\), is the price that the target receives. The light part, from
\(\max\{p,R\}\) to \(\theta\), is what the challenger keeps. Together the two parts cover \([0,\theta]\).}
\label{fig:auction-bars}
\end{figure}

\subsection*{The split: target and challenger share the whole value}

Figure~\ref{fig:auction-bars} suggests that the two amounts add up to \(\theta\). This holds at every incumbent
value.

\begin{derivation}[the split at a fixed incumbent value, for \(\theta\ge p\)]
Add the two amounts in each row of Table~\ref{tab:auction-cases}:
\begin{align*}
R<p\le\theta:\quad \max\{p,\min(R,\theta)\}+\bigl(\theta-\max\{p,R\}\bigr)_+&=p+(\theta-p)=\theta
\why{first row}\\
p\le R<\theta:\quad \max\{p,\min(R,\theta)\}+\bigl(\theta-\max\{p,R\}\bigr)_+&=R+(\theta-R)=\theta
\why{second row}\\
R\ge\theta:\quad \max\{p,\min(R,\theta)\}+\bigl(\theta-\max\{p,R\}\bigr)_+&=\theta+0=\theta
\why{third row; the tie}
\end{align*}
The sum equals \(\theta\) at every \(R\). Take the expectation over \(R\):
\begin{align*}
t(\theta)+g(\theta)&=\E_R\bigl[\max\{p,\min(R,\theta)\}\bigr]+\E_R\bigl[(\theta-\max\{p,R\})_+\bigr]
\why{definitions of \(t\) and \(g\)}\\
&=\E_R\bigl[\max\{p,\min(R,\theta)\}+(\theta-\max\{p,R\})_+\bigr] \why{linearity}\\
&=\E_R[\theta] \why{the sum is \(\theta\) at every \(R\)}\\
&=\theta \why{\(\theta\) is fixed}
\end{align*}
\end{derivation}

So, for every value at or above the reserve,
\begin{equation}\label{eq:auction-split}
t(\theta)+g(\theta)=\theta,\qquad \theta\ge p.
\end{equation}

\begin{intuition}[the split]
The second-price rule passes the whole value of an admissible challenger to the target and the challenger
together. The incumbent takes none of it: when the incumbent wins, it pays the full challenger value \(\theta\),
and its own gain \(R-\theta\) comes from its own value. The incumbent's value only decides how the challenger's
value is divided.
\end{intuition}

Below the reserve the split fails. For \(\theta<p\), the target receives \(p\,\1\{R\ge p\}\) and the challenger
keeps zero, so \(t(\theta)+g(\theta)=t_0\). An inadmissible value goes to no one.

\subsection*{Entry never lowers proceeds}

The paper notes that \(t_j\ge t_0\) for both classes. The reason: an admissible challenger raises the proceeds to
at least the reserve, and an inadmissible challenger changes nothing. Here is the argument at each fixed \(R\).

\begin{derivation}[\(t_j\ge t_0\)]
For \(\theta\ge p\) and every \(R\):
\begin{align*}
\max\{p,\min(R,\theta)\}&\ge p \why{a maximum is at least each entry}\\
&\ge p\,\1\{R\ge p\} \why{the indicator is at most one}
\end{align*}
For \(\theta<p\), the target receives \(p\,\1\{R\ge p\}\) with or without the challenger. So, at every \(R\),
the proceeds with the challenger are at least the proceeds without it. Take expectations:
\begin{align*}
t(\theta)&\ge\E_R\bigl[p\,\1\{R\ge p\}\bigr]=t_0 \why{expectations keep a weak inequality}\\
t_j=\E_j[t(\theta)]&\ge\E_j[t_0]=t_0 \why{\(t_0\) does not depend on \(\theta\)}
\end{align*}
\end{derivation}

Chapter~\ref{ch:trading} uses this fact: the entry gains \(t_H-t_0\) and \(t_L-t_0\) are never negative.

\begin{checkyourself}
\label{q:auction-jump}
Compute \(t(p)\) and \(g(p)\), the payoffs of a challenger whose value equals the reserve. Compare \(t(p)\) with
\(t_0\), the proceeds for any \(\theta<p\). Explain the difference in words: on which incumbent values do the two
cases give the target different amounts?
\end{checkyourself}

% ====================================================================
\section{The challenger's profit as an integral}\label{sec:auction-profit}

The goal of this section is the formula for \(g_j\). The method has three steps. First, write the challenger's
profit at a fixed \(R\) as a sum over levels \(y\). Then average over \(R\). Last, average over the class.

\subsection*{Step 1: the profit as a sum over levels}

Think of the challenger's value above the reserve as a stack of thin slices, one at each level \(y\) between \(p\)
and \(\theta\). At a fixed incumbent value, the challenger keeps the slice at level \(y\) exactly when the
incumbent's value is below \(y\). So, for \(\theta\ge p\),
\[
\bigl(\theta-\max\{p,R\}\bigr)_+=\int_p^\theta\1\{R<y\}\,dy.
\]
The integrand is one at the levels \(y\) above \(R\) and zero at the levels below \(R\). So the integral is the
length of the set of levels in \([p,\theta]\) that lie above \(R\). Check the three cases:
\begin{itemize}
\item \(R<p\): every level in \([p,\theta]\) lies above \(R\). The length is \(\theta-p\), and
\((\theta-\max\{p,R\})_+=\theta-p\).
\item \(p\le R<\theta\): the levels above \(R\) form \((R,\theta]\). The length is \(\theta-R\), and
\((\theta-\max\{p,R\})_+=\theta-R\).
\item \(R\ge\theta\): no level in \([p,\theta]\) lies above \(R\). The length is zero, and
\((\theta-\max\{p,R\})_+=0\).
\end{itemize}

\subsection*{Step 2: average over the incumbent}

The expectation of an integral of indicators is the integral of probabilities. The order of the two operations
can change because the integrand is never negative (Tonelli's theorem; Section~\ref{sec:toolkit-expectation}).

\begin{derivation}[the profit of a known value, \(\theta\ge p\)]
\begin{align*}
g(\theta)&=\E_R\bigl[(\theta-\max\{p,R\})_+\bigr] \why{definition of \(g\)}\\
&=\E_R\Bigl[\int_p^\theta\1\{R<y\}\,dy\Bigr] \why{Step 1}\\
&=\int_p^\theta\E_R\bigl[\1\{R<y\}\bigr]\,dy \why{change the order; integrand \(\ge0\)}\\
&=\int_p^\theta\Pr(R<y)\,dy \why{expectation of an indicator}\\
&=\int_p^\theta F(y)\,dy \why{\(F\) is continuous}
\end{align*}
\end{derivation}

This is formula~\eqref{eq:toolkit-fixed} of the tail formula section (Section~\ref{sec:toolkit-tail}).

\begin{intuition}[the profit of a known value]
The integral adds up, over the levels between the reserve and the challenger's value, the probability that the
incumbent is below each level. A level \(y\) pays the challenger when the incumbent's value lies under it.
\begin{itemize}
\item An incumbent that never meets the reserve has \(F(y)=1\) for \(y\ge p\). Then \(g(\theta)=\theta-p\): the
challenger always wins at the reserve.
\item An incumbent that is always above \(\theta\) has \(F(y)=0\) on \([p,\theta]\). Then \(g(\theta)=0\): the
challenger never wins.
\end{itemize}
\end{intuition}

For \(\theta<p\), the challenger does not bid, so \(g(\theta)=0\). One formula covers both cases:
\[
g(\theta)=\int_p^\infty\1\{y\le\theta\}\,F(y)\,dy\qquad\text{for every }\theta.
\]
For \(\theta\ge p\), the indicator keeps the levels in \([p,\theta]\), which gives \(\int_p^\theta F(y)\,dy\).
For \(\theta<p\), every level \(y\ge p\) lies above \(\theta\), so the indicator is zero and the integral is zero.

\subsection*{Step 3: average over the class}

The challenger does not know \(\theta\) when it decides to prepare. Its class-\(j\) profit is
\(g_j=\E_j[g(\theta)]\). This expectation is a double sum: over challenger values \(\theta\), and over levels
\(y\) with \(p\le y\le\theta\). Figure~\ref{fig:auction-region} shows the region of pairs \((\theta,y)\).

\begin{figure}[htbp]
\centering
\begin{tikzpicture}[x=0.8cm,y=0.8cm,every node/.style={font=\small}]
  \fill[accent!10] (1.5,1.5) -- (6,6) -- (7.2,6) -- (7.2,1.5) -- cycle;
  \draw[->] (0,0) -- (7.6,0) node[right] {\(\theta\)};
  \draw[->] (0,0) -- (0,6.4) node[above right] {level \(y\)};
  \draw[rulegrey,dashed] (0,1.5) -- (7.2,1.5);
  \draw[rulegrey,dashed] (0,0) -- (6,6);
  \node[anchor=east] at (0,1.5) {\(p\)};
  \draw (1.5,0.1) -- (1.5,-0.1) node[below] {\(p\)};
  \node[anchor=south] at (6,6) {\(y=\theta\)};
  % vertical strip at fixed theta
  \draw[accent,line width=1.6pt] (4.6,1.5) -- (4.6,4.6);
  \draw (4.6,0.1) -- (4.6,-0.1) node[below] {\(\theta\)};
  \node[anchor=west,align=left] at (4.7,3.9) {fix \(\theta\):\\ \(y\) runs from \(p\) to \(\theta\)};
  % horizontal strip at fixed y
  \draw[notegrey,line width=1.6pt] (2.8,2.8) -- (7.2,2.8);
  \node[anchor=east] at (0,2.8) {\(y\)};
  \draw[rulegrey,dotted] (0,2.8) -- (2.8,2.8);
  \node[anchor=north west,align=left] at (5.2,2.7) {fix \(y\):\\ \(\theta\) runs from \(y\) up};
\end{tikzpicture}
\caption{The region of the double sum in \(g_j\): all pairs of a challenger value \(\theta\) and a level \(y\)
with \(p\le y\le\theta\). The first order of
summation fixes \(\theta\) and adds over the levels \(y\) in a vertical strip; this gives \(g(\theta)\). The
second order fixes a level \(y\) and collects the probability of every challenger value \(\theta\ge y\) in a
horizontal strip; this gives \(\Pr_j(\theta\ge y)\).}
\label{fig:auction-region}
\end{figure}

The two orders give the same total, because every term is nonnegative (Tonelli's theorem again). The second order
is the useful one: at each level \(y\), it asks how likely the challenger's value is to reach \(y\).

\begin{derivation}[the class-\(j\) gross profit]
\begin{align*}
g_j&=\E_j\bigl[g(\theta)\bigr] \why{definition of \(g_j\)}\\
&=\E_j\Bigl[\int_p^\infty\1\{y\le\theta\}\,F(y)\,dy\Bigr] \why{one formula for every \(\theta\)}\\
&=\int_p^\infty\E_j\bigl[\1\{y\le\theta\}\bigr]\,F(y)\,dy \why{change the order; terms \(\ge0\)}\\
&=\int_p^\infty\Pr\nolimits_j(\theta\ge y)\,F(y)\,dy \why{expectation of an indicator}\\
&=\int_p^\infty\bigl[1-G_j(y^-)\bigr]\,F(y)\,dy \why{\(\Pr_j(\theta\ge y)=1-G_j(y^-)\)}\\
&=\int_p^\infty\bigl[1-G_j(y)\bigr]\,F(y)\,dy \why{the left limit differs at atoms only}
\end{align*}
\end{derivation}

The last step needs a word. The functions \(G_j(y^-)\) and \(G_j(y)\) differ only at the atoms of \(G_j\). A
distribution function has at most countably many atoms. At most one atom has a mass above \(\tfrac12\), at most
two have a mass above \(\tfrac13\), and so on, so the atoms can be listed one by one. Each single point has zero
length. An integral over \(y\) measures area, and a countable list of single points adds no area. So the change at
the atoms does not change the integral.

The result is the second line of the paper's equation for Proposition~1:
\begin{equation}\label{eq:auction-gross}
g_j(F)=\int_p^\infty\bigl[1-G_j(y)\bigr]\,F(y)\,dy,\qquad j\in\{H,L\}.
\end{equation}

\begin{intuition}[the gross profit formula]
A level \(y\) above the reserve earns for the challenger when two events happen together. The challenger's value
lies above \(y\), with probability \(1-G_j(y)\). The incumbent's value lies below \(y\), with probability
\(F(y)\). The two values are independent, so the probability of both events is the product. The gross profit
adds these products over all levels above the reserve.
\end{intuition}

The integral is finite. The factor \(F(y)\) lies between zero and one, so
\[
0\le g_j\le\int_p^\infty\bigl[1-G_j(y)\bigr]\,dy\le\int_0^\infty\bigl[1-G_j(y)\bigr]\,dy=\E_j[\theta]<\infty,
\]
where the equality is the tail formula for a nonnegative value (Section~\ref{sec:toolkit-tail}).

The incumbent enters \eqref{eq:auction-gross} only through \(F\), and \(F\) carries a nonnegative weight. A
strengthening lowers \(F\), so it cannot raise \(g_j\). Section~\ref{sec:auction-prop1} makes this exact.

\begin{checkyourself}
\label{q:auction-pointgross}
Take point-mass classes: a class-\(H\) value equals \(h\) and a class-\(L\) value equals \(\ell\), with
\(p<\ell<h\). Use \eqref{eq:auction-gross} to show that \(g_H-g_L=\int_\ell^hF(y)\,dy\). Say in words which
incumbent values this integral counts, and why \(g_H\ge g_L\).
\end{checkyourself}

% ====================================================================
\section{Target proceeds and the spread}\label{sec:auction-spread}

This section follows the paper's proof for target proceeds line by line. The steps mirror the steps for the
profit, with one extra term from the reserve.

\subsection*{Step 1: the proceeds as a sum over levels}

For \(\theta\ge p\), the target receives the reserve plus every slice between \(p\) and \(\theta\) that lies below
the incumbent's value:
\[
\max\{p,\min(R,\theta)\}=p+\int_p^\theta\1\{R>y\}\,dy.
\]
The integral is the length of the set of levels in \([p,\theta]\) that lie below \(R\). Check the three cases:
\begin{itemize}
\item \(R<p\): no level in \([p,\theta]\) lies below \(R\). The right side is \(p+0=p\), the first row of
Table~\ref{tab:auction-cases}.
\item \(p\le R<\theta\): the levels below \(R\) form \([p,R)\), of length \(R-p\). The right side is
\(p+(R-p)=R\).
\item \(R\ge\theta\): every level in \([p,\theta]\) lies below \(R\), except the single point \(y=R\) in the tie.
The length is \(\theta-p\), and the right side is \(\theta\).
\end{itemize}

\subsection*{Step 2: average over the incumbent}

\begin{derivation}[the proceeds for a known value]
For \(\theta\ge p\):
\begin{align*}
t(\theta)&=\E_R\bigl[\max\{p,\min(R,\theta)\}\bigr] \why{definition of \(t\)}\\
&=p+\E_R\Bigl[\int_p^\theta\1\{R>y\}\,dy\Bigr] \why{Step 1; linearity}\\
&=p+\int_p^\theta\Pr(R>y)\,dy \why{change the order; indicator}\\
&=p+\int_p^\theta\bigl[1-F(y)\bigr]\,dy \why{\(F\) is continuous}
\end{align*}
For \(\theta<p\):
\begin{align*}
t(\theta)&=\E_R\bigl[p\,\1\{R\ge p\}\bigr] \why{the challenger does not bid}\\
&=p\,\Pr(R\ge p)=p\,\bigl[1-F(p)\bigr]=t_0 \why{as in Section~\ref{sec:auction-setup}}
\end{align*}
\end{derivation}

As a check, add \(g(\theta)=\int_p^\theta F(y)\,dy\) to \(t(\theta)\) for \(\theta\ge p\). The paper puts the
same check in words: the indicators \(\1\{R>y\}\) and \(\1\{R<y\}\) add up to one at every level except the single
point \(y=R\), which has zero length.
\begin{align*}
t(\theta)+g(\theta)&=p+\int_p^\theta\bigl[1-F(y)\bigr]\,dy+\int_p^\theta F(y)\,dy \why{the two formulas}\\
&=p+\int_p^\theta1\,dy \why{add the integrands}\\
&=\theta \why{the integral is \(\theta-p\)}
\end{align*}
This is the split \eqref{eq:auction-split} again.

\subsection*{Step 3: the jump at the reserve}

At \(\theta=p\), the formula for \(\theta\ge p\) gives \(t(p)=p+0=p\). Just below the reserve,
\(t(\theta)=t_0=p-pF(p)\). So \(t\) jumps up by \(pF(p)\) at \(\theta=p\). The reason: when the incumbent fails
the reserve, which has probability \(F(p)\), an admissible challenger buys at \(p\), and an inadmissible one
leaves the target unsold.

With the jump written out, one formula holds for every \(\theta\):
\[
t(\theta)=t_0+pF(p)\,\1\{\theta\ge p\}+\int_p^\infty\1\{y\le\theta\}\bigl[1-F(y)\bigr]\,dy.
\]

\begin{derivation}[the formula for every \(\theta\)]
For \(\theta\ge p\):
\begin{align*}
&t_0+pF(p)+\int_p^\infty\1\{y\le\theta\}\bigl[1-F(y)\bigr]\,dy\\
&\qquad=p\bigl[1-F(p)\bigr]+pF(p)+\int_p^\theta\bigl[1-F(y)\bigr]\,dy
\why{\(t_0=p[1-F(p)]\); the indicator}\\
&\qquad=p+\int_p^\theta\bigl[1-F(y)\bigr]\,dy=t(\theta) \why{Step 2}
\end{align*}
For \(\theta<p\), both the indicator \(\1\{\theta\ge p\}\) and every \(\1\{y\le\theta\}\) with \(y\ge p\) are
zero, so the right side is \(t_0=t(\theta)\).
\end{derivation}

\subsection*{Step 4: average over the class}

Two facts about the class distribution enter: \(\E_j[\1\{\theta\ge p\}]=\Pr_j(\theta\ge p)=1-G_j(p^-)\) and
\(\E_j[\1\{y\le\theta\}]=\Pr_j(\theta\ge y)=1-G_j(y^-)\).

\begin{derivation}[the class-\(j\) proceeds]
\begin{align*}
t_j&=\E_j\bigl[t(\theta)\bigr] \why{definition of \(t_j\)}\\
&=t_0+pF(p)\,\E_j\bigl[\1\{\theta\ge p\}\bigr]
+\E_j\Bigl[\int_p^\infty\1\{y\le\theta\}\bigl[1-F(y)\bigr]\,dy\Bigr] \why{Step 3; linearity}\\
&=t_0+pF(p)\,\E_j\bigl[\1\{\theta\ge p\}\bigr]
+\int_p^\infty\E_j\bigl[\1\{y\le\theta\}\bigr]\bigl[1-F(y)\bigr]\,dy \why{change the order}\\
&=t_0+pF(p)\bigl[1-G_j(p^-)\bigr]+\int_p^\infty\bigl[1-G_j(y^-)\bigr]\bigl[1-F(y)\bigr]\,dy
\why{the two facts}\\
&=t_0+pF(p)\bigl[1-G_j(p^-)\bigr]+\int_p^\infty\bigl[1-G_j(y)\bigr]\bigl[1-F(y)\bigr]\,dy
\why{atoms add no area}
\end{align*}
\end{derivation}

\begin{watchout}
The last step replaces \(G_j(y^-)\) by \(G_j(y)\) inside the integral over \(y\). The left limit \(G_j(p^-)\) in
the middle term stays. An atom at \(p\) makes \(G_j(p^-)\) and \(G_j(p)\) differ. The middle term is a single
number, so it must keep the left limit. The integral does not notice the difference at one point.
\end{watchout}

\subsection*{Step 5: the difference of the classes}

\begin{derivation}[the spread]
\begin{align*}
\DT&=t_H-t_L \why{definition}\\
&=pF(p)\Bigl\{\bigl[1-G_H(p^-)\bigr]-\bigl[1-G_L(p^-)\bigr]\Bigr\}\\
&\qquad+\int_p^\infty\Bigl\{\bigl[1-G_H(y)\bigr]-\bigl[1-G_L(y)\bigr]\Bigr\}\bigl[1-F(y)\bigr]\,dy
\why{Step 4; \(t_0\) cancels}\\
&=pF(p)\bigl[G_L(p^-)-G_H(p^-)\bigr]+\int_p^\infty\bigl[G_L(y)-G_H(y)\bigr]\bigl[1-F(y)\bigr]\,dy
\why{the ones cancel}\\
&=pF(p)\,D(p^-)+\int_p^\infty\bigl[1-F(y)\bigr]\,D(y)\,dy \why{definition of \(D\)}
\end{align*}
Here \(D(p^-)=G_L(p^-)-G_H(p^-)\) is the left limit of \(D\) at \(p\).
\end{derivation}

This is the first line of the paper's equation for Proposition~1:
\begin{equation}\label{eq:auction-spread}
\DT(F)=pF(p)\,D(p^-)+\int_p^\infty\bigl[1-F(y)\bigr]\,D(y)\,dy.
\end{equation}

Each integral is finite: \(0\le\int_p^\infty[1-G_j(y)]\,dy\le\E_j[\theta]<\infty\), as in
Section~\ref{sec:auction-profit}, and \(0\le F\le1\). So the difference of the two class integrals is a
difference of two finite numbers.

\subsection*{What the two terms say}

The integral term adds up the levels above the reserve. At the level \(y\), the factor \(1-F(y)\) is the
probability that the incumbent's value lies above \(y\). On that event, the target collects the slice at \(y\)
from any challenger whose value passes \(y\). The factor \(D(y)\) says how much more likely a high-class value is
to pass \(y\). So \([1-F(y)]D(y)\,dy\) is the class difference in the slice that reaches the target.

The reserve term \(pF(p)D(p^-)\) has three factors:
\begin{itemize}
\item \(F(p)\) is the probability that the incumbent fails the reserve.
\item On that event, the target receives \(p\) from an admissible challenger and nothing from an inadmissible one.
\item The class difference in the chance to be admissible is
\(\bigl[1-G_H(p^-)\bigr]-\bigl[1-G_L(p^-)\bigr]=G_L(p^-)-G_H(p^-)=D(p^-)\).
\end{itemize}
Under the class ranking~\eqref{eq:auction-ranking}, \(G_H(p^-)=G_L(p^-)=0\), so \(D(p^-)=0\) and the reserve term
is zero. The term is also zero under equal admissibility, and when the incumbent always meets the reserve,
\(F(p)=0\).

\begin{watchout}
Without equal admissibility, the reserve term can work against the rest of the formula. Suppose the low class
fails the reserve more often, so \(D(p^-)>0\). A stronger incumbent fails the reserve less often,
\(\tilde F(p)\le F(p)\), so the term \(pF(p)D(p^-)\) shrinks. The paper notes that this loss can outweigh the
slices that the strengthening moves to the target. The class ranking closes this channel. When the classes are
ranked above the reserve, this term is the only way in which a strengthening can narrow the spread
(Section~\ref{sec:auction-weakest}).
\end{watchout}

\subsection*{Two signs under the class ranking}

Under the class ranking, \(D(y)\ge0\) for every \(y\) and \(D(p^-)=0\). Two signs follow.
\begin{align*}
\DT&=\int_p^\infty\bigl[1-F(y)\bigr]\,D(y)\,dy\ \ge0 \why{\(D(p^-)=0\); both factors \(\ge0\)}\\
g_H-g_L&=\int_p^\infty\bigl[1-G_H(y)\bigr]F(y)\,dy-\int_p^\infty\bigl[1-G_L(y)\bigr]F(y)\,dy
\why{\eqref{eq:auction-gross}}\\
&=\int_p^\infty\bigl[G_L(y)-G_H(y)\bigr]F(y)\,dy \why{combine the integrands}\\
&=\int_p^\infty D(y)\,F(y)\,dy\ \ge0 \why{both factors \(\ge0\)}
\end{align*}
So the high class gives the target at least as much, and it earns at least as much profit.

\emph{Extremes.} If the incumbent never meets the reserve, \(F(y)=1\) for \(y\ge p\), and the integral term
vanishes: \(\DT=pD(p^-)\), which is zero under the class ranking. Next, suppose the incumbent is always above the
reserve and above every challenger value. Then \(F(p)=0\), and \(1-F(y)=1\) wherever \(D(y)\ne0\). So
\(\DT=\int_p^\infty D(y)\,dy\), the whole class value gap above the reserve.

% ====================================================================
\section{One incumbent value: the slices}\label{sec:auction-slices}

The formula \eqref{eq:auction-spread} averages over the incumbent. The paper also reads it at one fixed incumbent
value. This view explains the word \emph{slice}, and it drives part~(c) in Section~\ref{sec:auction-weakest}.

\subsection*{The target's proceeds at one incumbent value}

Fix an incumbent value \(R\ge p\). Then the target receives
\[
\min\bigl\{R,\max\{\theta,p\}\bigr\}\qquad\text{for every challenger value }\theta.
\]
The reason is that an admissible challenger and the incumbent both bid at least \(p\), and an inadmissible
challenger leaves the incumbent to pay \(p\). Check the cases:
\begin{itemize}
\item \(\theta<p\): the incumbent is the only admissible bidder and pays \(p\). The formula gives
\(\min\{R,p\}=p\), since \(R\ge p\).
\item \(p\le\theta\le R\): the incumbent wins, or ties, and the price is \(\theta\). The formula gives
\(\min\{R,\theta\}=\theta\).
\item \(\theta>R\): the challenger wins and pays \(\max\{p,R\}=R\). The formula gives \(\min\{R,\theta\}=R\).
\end{itemize}
If instead \(R<p\), the incumbent does not bid, and the target receives \(p\,\1\{\theta\ge p\}\).

For \(R\ge p\), write the proceeds as the reserve plus the slices between \(p\) and \(R\) that the challenger's
value passes:
\[
\min\bigl\{R,\max\{\theta,p\}\bigr\}=p+\int_p^R\1\{\theta>y\}\,dy.
\]
The integral is the length of the levels in \([p,R]\) below \(\theta\). It is zero when \(\theta\le p\), it is
\(\theta-p\) when \(p<\theta\le R\), and it is \(R-p\) when \(\theta>R\). In each case the right side matches the
three cases above.

\begin{derivation}[the expected proceeds of class \(j\) at one incumbent value \(R\ge p\)]
\begin{align*}
\E_j\Bigl[\min\bigl\{R,\max\{\theta,p\}\bigr\}\Bigr]&=p+\E_j\Bigl[\int_p^R\1\{\theta>y\}\,dy\Bigr]
\why{the slice formula; linearity}\\
&=p+\int_p^R\Pr\nolimits_j(\theta>y)\,dy \why{change the order; indicator}\\
&=p+\int_p^R\bigl[1-G_j(y)\bigr]\,dy \why{\(\Pr_j(\theta>y)=1-G_j(y)\)}
\end{align*}
\end{derivation}

Notice the strict inequality \(\theta>y\) here. Its probability is \(1-G_j(y)\) exactly, so no left limit appears.
In Section~\ref{sec:auction-profit} the indicator was \(\1\{y\le\theta\}\), with the probability
\(\Pr_j(\theta\ge y)=1-G_j(y^-)\). Both routes lead to the same integrals, since atoms add no area.

\subsection*{The class difference at one incumbent value}

The paper calls the expected class difference in proceeds at the incumbent value \(R\) the function \(\Gamma\):
\[
\Gamma(R)=
\begin{cases}
pD(p^-), & R<p,\\[3pt]
\int_p^RD(y)\,dy, & R\ge p.
\end{cases}
\]

\begin{derivation}[the two cases of \(\Gamma\)]
For \(R\ge p\):
\begin{align*}
\Gamma(R)&=\Bigl(p+\int_p^R\bigl[1-G_H(y)\bigr]dy\Bigr)-\Bigl(p+\int_p^R\bigl[1-G_L(y)\bigr]dy\Bigr)
\why{the last derivation, twice}\\
&=\int_p^R\bigl[G_L(y)-G_H(y)\bigr]\,dy \why{\(p\) cancels; combine}\\
&=\int_p^RD(y)\,dy \why{definition of \(D\)}
\end{align*}
For \(R<p\), the target receives \(p\,\1\{\theta\ge p\}\):
\begin{align*}
\Gamma(R)&=p\,\Pr\nolimits_H(\theta\ge p)-p\,\Pr\nolimits_L(\theta\ge p) \why{expectation of an indicator}\\
&=p\bigl[1-G_H(p^-)\bigr]-p\bigl[1-G_L(p^-)\bigr] \why{left limits}\\
&=p\,D(p^-) \why{definition of \(D(p^-)\)}
\end{align*}
\end{derivation}

The spread is the average of \(\Gamma\) over the incumbent's value:
\[
\DT(F)=\E_F\bigl[\Gamma(R)\bigr].
\]
The reason: \(t_j\) averages the proceeds over \(\theta\) and over \(R\). The two values are independent and the
proceeds are never negative, so the order of the two averages does not matter (Tonelli's theorem). Average over
\(\theta\) first, at each fixed \(R\), and then over \(R\). The class difference of the inner averages is
\(\Gamma(R)\).

\begin{derivation}[{\(\E_F[\Gamma(R)]\) gives back \eqref{eq:auction-spread}}]
For every \(R\), \(\1\{R\ge p\}\int_p^RD(y)\,dy=\int_p^\infty\1\{y<R\}D(y)\,dy\), since both sides are zero for
\(R<p\). Write \(D=[1-G_H]-[1-G_L]\), a difference of two nonnegative functions, and change the order for each
part separately.
\begin{align*}
\E_F\bigl[\Gamma(R)\bigr]&=\Pr(R<p)\,pD(p^-)+\E_F\Bigl[\int_p^\infty\1\{y<R\}\,D(y)\,dy\Bigr]
\why{the two cases of \(\Gamma\)}\\
&=F(p)\,pD(p^-)+\int_p^\infty\Pr(R>y)\,D(y)\,dy \why{change the order, part by part}\\
&=pF(p)\,D(p^-)+\int_p^\infty\bigl[1-F(y)\bigr]\,D(y)\,dy \why{\(F\) is continuous}
\end{align*}
\end{derivation}

\subsection*{The slices}

The class value gap above the reserve is the sum of thin slices:
\[
\E_H\bigl[\max\{\theta,p\}\bigr]-\E_L\bigl[\max\{\theta,p\}\bigr]=\int_p^\infty D(y)\,dy.
\]
This follows from the tail formula; Section~\ref{sec:auction-gap} derives it step by step.

The slice at level \(y\), of thickness \(dy\), is present in the challenger's value when \(\theta>y\). It is
present with probability \(1-G_H(y)\) in class \(H\) and with probability \(1-G_L(y)\) in class \(L\). So
\(D(y)\,dy\) is the class difference in this slice: one slice of the class value gap.

At the incumbent value \(R\ge p\), the gap splits at \(R\):
\begin{itemize}
\item The slices below \(R\) reach the target. Their sum is \(\Gamma(R)=\int_p^RD(y)\,dy\).
\item The slices above \(R\) stay with the challenger. The challenger keeps \((\theta-R)_+=\int_R^\infty\1\{\theta>y\}\,dy\),
so the class difference in its expected profit is \(\int_R^\infty D(y)\,dy\).
\end{itemize}
In the paper's words, each slice \(D(y)\,dy\) at a level \(y>p\) reaches the target when the incumbent outbids the
challenger at that level and pays its bid. Otherwise the slice stays with the challenger. Use the clock picture of
Section~\ref{sec:auction-story}. Suppose that the slice is present, \(\theta>y\), and that the incumbent is still
in at level \(y\), \(R>y\). Then the price passes \(y\): either the incumbent wins and pays \(\theta\), or the
challenger wins and pays \(R\). Figure~\ref{fig:auction-band}
draws the split.

\begin{figure}[htbp]
\centering
\begin{tikzpicture}[x=1.25cm,y=3cm,every node/.style={font=\small}]
  \def\pp{1}
  \def\RR{4}
  \fill[accent!40] plot[domain=\pp:\RR,samples=60] (\x,{exp(-((\x-\pp)/5)^2)})
     -- plot[domain=\RR:\pp,samples=60] (\x,{exp(-((\x-\pp)/2.6)^2)}) -- cycle;
  \fill[notegrey!22] plot[domain=\RR:9.5,samples=60] (\x,{exp(-((\x-\pp)/5)^2)})
     -- plot[domain=9.5:\RR,samples=60] (\x,{exp(-((\x-\pp)/2.6)^2)}) -- cycle;
  \draw[->] (0,0) -- (10,0) node[below left] {level \(y\)};
  \draw[->] (0,0) -- (0,1.2);
  \draw (0.08,1) -- (-0.08,1) node[left] {\(1\)};
  \draw[ink,thick] plot[domain=\pp:9.5,samples=80] (\x,{exp(-((\x-\pp)/5)^2)});
  \draw[ink,thick,dashed] plot[domain=\pp:9.5,samples=80] (\x,{exp(-((\x-\pp)/2.6)^2)});
  \node[anchor=south west] at (6.6,0.36) {\(1-G_H(y)\)};
  \node[anchor=north east] at (3.0,0.45) {\(1-G_L(y)\)};
  \draw (\pp,0.03) -- (\pp,-0.03) node[below] {\(p\)};
  \draw[accent,thick] (\RR,0) -- (\RR,1.05) node[above] {\(R\)};
  \node[accent,anchor=south] at (2.3,1.02) {to the target};
  \draw[accent] (2.3,1.02) -- (3.4,0.62);
  \node[notegrey,anchor=west,align=left] at (7.3,0.9) {kept by the\\ challenger};
  \draw[notegrey] (7.3,0.85) -- (5.6,0.5);
\end{tikzpicture}
\caption{The class value gap as the area between the two curves \(1-G_H(y)\) and \(1-G_L(y)\) above the reserve.
The height of the band at level \(y\) is \(D(y)\). At an incumbent value \(R\), the band to the left of \(R\)
reaches the target; its area is \(\Gamma(R)\). The band to the right of \(R\) stays with the challenger. A higher
\(R\) moves the boundary to the right. The curves are a sketch under the class ranking.}
\label{fig:auction-band}
\end{figure}

\begin{intuition}[why a stronger incumbent widens the spread]
A stronger incumbent makes \(R>y\) more likely at every level \(y\). So, in expectation, more slices reach the
target. Under the class ranking every slice \(D(y)\,dy\) is nonnegative. A move of nonnegative slices to the
target cannot lower the class difference in proceeds.
\end{intuition}

\subsection*{An illustration with point-mass classes}

Take the paper's point-mass classes: a class-\(H\) value equals \(h\), a class-\(L\) value equals \(\ell\), and
\(p<\ell<h\). Then \(G_H(y)=\1\{y\ge h\}\) and \(G_L(y)=\1\{y\ge\ell\}\), so
\[
D(y)=\1\{y\ge\ell\}-\1\{y\ge h\}=\1\{\ell\le y<h\},\qquad D(p^-)=0.
\]
For \(R\ge p\), \(\Gamma(R)\) is the length of \([p,R]\cap[\ell,h)\). So \(\Gamma(R)=0\) for \(R\le\ell\),
\(\Gamma(R)=R-\ell\) for \(\ell\le R<h\), and \(\Gamma(R)=h-\ell\) for \(R\ge h\). The paper writes
\(\Gamma(R)=(R-\ell)_+\) for \(R<h\). Panel~(a) of Figure~\ref{fig:auction-kernel} draws this shape.

Each case has a direct explanation:
\begin{itemize}
\item \(R<\ell\): both classes beat the incumbent and pay the same price \(\max\{p,R\}\). The class makes no
difference to the target.
\item \(\ell\le R<h\): the high class wins and pays \(R\). The low class loses, and the incumbent pays \(\ell\).
The difference is \(R-\ell\).
\item \(R\ge h\): both classes lose, and the incumbent pays \(h\) or \(\ell\). The difference is the whole gap
\(h-\ell\).
\end{itemize}
Averaged over the incumbent, \eqref{eq:auction-spread} gives \(\DT=\int_\ell^h[1-F(y)]\,dy\). When \(\bar r<h\),
the factor \(1-F(y)\) is zero above \(\bar r\), so \(\DT=\int_\ell^{\bar r}[1-F(y)]\,dy\). The same classes give
\(g_j=\int_p^\theta F(y)\,dy\) at \(\theta=h\) and \(\theta=\ell\). Chapter~\ref{ch:example} computes these
integrals for the uniform incumbent of the example.

\begin{figure}[htbp]
\centering
\begin{tikzpicture}[x=0.7cm,y=0.85cm,every node/.style={font=\small}]
  % panel (a): p < ell < h
  \begin{scope}
    \draw[->] (0,0) -- (7.6,0) node[below] {\(R\)};
    \draw[->] (0,0) -- (0,3.2) node[left] {\(\Gamma(R)\)};
    \draw[accent,very thick] (0,0) -- (3,0) -- (5.5,2.5) -- (7.3,2.5);
    \foreach \x/\t in {1/p,3/\ell,5.5/h} {\draw (\x,0.08) -- (\x,-0.08) node[below] {\(\t\)};}
    \draw[rulegrey,dashed] (0,2.5) -- (5.5,2.5);
    \node[left] at (0,2.5) {\(h-\ell\)};
    \node at (3.8,-1.2) {(a) \(p<\ell<h\)};
  \end{scope}
  % panel (b): ell < p < h
  \begin{scope}[xshift=7.9cm]
    \draw[->] (0,0) -- (7.6,0) node[below] {\(R\)};
    \draw[->] (0,0) -- (0,3.2) node[left] {\(\Gamma(R)\)};
    \draw[accent,very thick] (0,1.6) -- (2.4,1.6);
    \draw[accent,very thick] (2.4,0) -- (5.5,1.9) -- (7.3,1.9);
    \draw[accent,fill=white] (2.4,1.6) circle (2.2pt);
    \fill[accent] (2.4,0) circle (2.2pt);
    \foreach \x/\t in {1/\ell,2.4/p,5.5/h} {\draw (\x,0.08) -- (\x,-0.08) node[below] {\(\t\)};}
    \node[left] at (0,1.6) {\(p\)};
    \draw[rulegrey,dashed] (2.4,1.9) -- (5.5,1.9);
    \node[right] at (7.3,1.9) {\(h-p\)};
    \node at (3.8,-1.2) {(b) \(\ell<p<h\)};
  \end{scope}
\end{tikzpicture}
\caption{The class difference in proceeds \(\Gamma(R)\) at a fixed incumbent value, for point-mass classes. In
panel~(a) the reserve lies below both values; \(\Gamma\) never falls, so every strengthening widens the spread or
leaves it unchanged. In panel~(b) the reserve excludes the low class; \(\Gamma\) drops from \(p\) to zero where the
incumbent starts to meet the reserve. The relative heights in panel~(b) are a sketch.}
\label{fig:auction-kernel}
\end{figure}

\begin{checkyourself}
\label{q:auction-slices}
In the point-mass case with \(p<\ell<h\), fix an incumbent value \(R\) with \(\ell\le R<h\). Find the
challenger's profit in each class at this \(R\). Show that \(\Gamma(R)\) plus the class difference in the
challenger's profit equals \(h-\ell\), and say which slices go where.
\end{checkyourself}

% ====================================================================
\section{Proposition 1: competition and the two returns to information}\label{sec:auction-prop1}

All the pieces are in place. Here is the paper's statement.

\needspace{10\baselineskip}
\begin{paperresult}{Proposition 1 (Competition and the two returns to information)}
Let the incumbent's value have a continuous distribution \(F\) on \([0,\bar r]\), extended by unity above
\(\bar r\). For any classes \(G_H,G_L\) with finite means,
\[
\begin{aligned}
\DT(F)&=pF(p)D(p^-)+\int_p^\infty[1-F(y)]\,D(y)\,dy,\\
g_j(F)&=\int_p^\infty[1-G_j(y)]\,F(y)\,dy,\qquad j\in\{H,L\}.
\end{aligned}
\]
Let \(\tilde F\le F\) be a continuous first-order stochastic strengthening of \(F\).
\begin{enumerate}[label=(\alph*),leftmargin=2em]
\item For every pair of classes, the strengthening weakly decreases \(g_H\) and \(g_L\), and hence gross
acquisition profit at every fixed belief.
\item Under the class ranking~\eqref{eq:auction-ranking}, the strengthening weakly increases the target-payoff
spread.
\end{enumerate}
The decrease in \(g_j\) is strict when \(\int_p^\infty[1-G_j(y)][F(y)-\tilde F(y)]\,dy>0\). The increase in (b) is
strict when \(\int_p^\infty[F(y)-\tilde F(y)]\,D(y)\,dy>0\).
\end{paperresult}

In plain words, line by line:
\begin{itemize}
\item \emph{Hypotheses.} \(F\) and \(\tilde F\) are continuous, which makes ties null and gives
\(\Pr(R<y)=F(y)\). The classes are any two distributions with finite means, atoms included, so every integral is
finite.
\item \emph{The two formulas} are \eqref{eq:auction-spread} and \eqref{eq:auction-gross}. Neither needs the class
ranking.
\item \emph{Part (a).} Each class profit falls or stays the same, for every pair of classes, and so does
\(B(\mu)\) at every fixed belief \(\mu\).
\item \emph{Part (b).} Under the class ranking, the spread rises or stays the same. Each change is strict when the
stated integral is positive.
\end{itemize}

\begin{derivation}[proof of Proposition~1]
The two formulas are \eqref{eq:auction-gross} and \eqref{eq:auction-spread}, derived in
Sections~\ref{sec:auction-profit} and~\ref{sec:auction-spread} for any classes with finite means.
\begin{enumerate}[label=\textit{Step \arabic*.},leftmargin=*,itemsep=0.6em]
\item \emph{The change in \(g_j\).} Apply \eqref{eq:auction-gross} at \(F\) and at \(\tilde F\):
\begin{align*}
g_j(F)-g_j(\tilde F)&=\int_p^\infty\bigl[1-G_j(y)\bigr]F(y)\,dy\\
&\qquad-\int_p^\infty\bigl[1-G_j(y)\bigr]\tilde F(y)\,dy
\why{\eqref{eq:auction-gross} twice}\\
&=\int_p^\infty\bigl[1-G_j(y)\bigr]\bigl[F(y)-\tilde F(y)\bigr]\,dy \why{combine the integrands}\\
&\ge0 \why{\(1-G_j\ge0\) and \(F-\tilde F\ge0\)}
\end{align*}
The weight \(1-G_j(y)\) is a probability, so it is nonnegative. The factor \(F(y)-\tilde F(y)\) is nonnegative
because \(\tilde F\le F\). So \(g_j\) falls weakly. The fall equals the integral, so it is strict exactly when
the integral is positive. No property of the classes enters this step.
\item \emph{Gross profit at a fixed belief.} At a fixed \(\mu\in[0,1]\):
\begin{align*}
B_F(\mu)-B_{\tilde F}(\mu)&=(1-\mu)\bigl[g_L(F)-g_L(\tilde F)\bigr]+\mu\bigl[g_H(F)-g_H(\tilde F)\bigr]
\why{\(B=(1-\mu)g_L+\mu g_H\)}\\
&\ge0 \why{Step 1; weights \(\ge0\)}
\end{align*}
The weights \(1-\mu\) and \(\mu\) are fixed and nonnegative, so an average of two weak falls is a weak fall.
\item \emph{The change in \(\DT\).} Apply \eqref{eq:auction-spread} at \(\tilde F\) and at \(F\):
\begin{align*}
\DT(\tilde F)-\DT(F)&=p\tilde F(p)D(p^-)+\int_p^\infty\bigl[1-\tilde F(y)\bigr]D(y)\,dy\\
&\qquad-pF(p)D(p^-)-\int_p^\infty\bigl[1-F(y)\bigr]D(y)\,dy \why{\eqref{eq:auction-spread} twice}\\
&=\int_p^\infty\bigl[F(y)-\tilde F(y)\bigr]D(y)\,dy-p\bigl[F(p)-\tilde F(p)\bigr]D(p^-)
\why{collect terms}
\end{align*}
In the integrals, \([1-\tilde F(y)]-[1-F(y)]=F(y)-\tilde F(y)\). In the reserve terms,
\(\tilde F(p)-F(p)=-[F(p)-\tilde F(p)]\). This difference formula holds for any classes.
\item \emph{The sign under the class ranking.} The class ranking gives \(D(p^-)=0\), so the second term is zero.
It also gives \(D(y)\ge0\), and \(F(y)-\tilde F(y)\ge0\) holds for every \(y\). So the integrand is nonnegative:
\begin{align*}
\DT(\tilde F)-\DT(F)&=\int_p^\infty\bigl[F(y)-\tilde F(y)\bigr]D(y)\,dy \why{\(D(p^-)=0\)}\\
&\ge0 \why{both factors \(\ge0\)}
\end{align*}
The rise equals the integral, so it is strict exactly when the integral is positive.
\item \emph{A weaker hypothesis.} Equal admissibility also makes \(D(p^-)=0\). Together with \(D(y)\ge0\) for
\(y>p\), it gives the same conclusion, since a single point \(y=p\) adds no area.
\end{enumerate}
\end{derivation}

\begin{intuition}[Proposition~1]
Part (a): in \eqref{eq:auction-gross}, each level pays the challenger only when the incumbent is below it. A
stronger incumbent is below each level less often. Part (b): in \eqref{eq:auction-spread}, each slice of the
class gap reaches the target only when the incumbent is above its level. A stronger incumbent is above each level
more often, and under the class ranking every slice is nonnegative. The same change in \(F\) works against the
challenger in one formula and for the target in the other.
\end{intuition}

\begin{watchout}
Proposition~1 gives weak inequalities. A strengthening can leave a payoff unchanged.
\begin{itemize}
\item If \(\tilde F\) and \(F\) differ only below the reserve, then \(F(y)=\tilde F(y)\) for \(y\ge p\), and
neither \(g_j\) nor the integral in (b) changes. An incumbent below the reserve never bids, so a move among such
values changes nothing for the challenger.
\item If the strengthening changes \(F\) only at levels above every class-\(j\) value, then \(1-G_j(y)=0\) there,
and \(g_j\) does not change. The incumbent beats class \(j\) at those levels anyway.
\item In the point-mass case, \(D(y)=0\) outside \([\ell,h)\). A strengthening that leaves \(F\) unchanged on
\([\ell,h)\) leaves the spread unchanged.
\end{itemize}
A strict change needs the stated integral to be positive.
\end{watchout}

\begin{watchout}
Proposition~1 compares payoffs at two incumbent distributions, with the classes and the reserve fixed. It says
nothing yet about equilibrium entry. A comparison of equilibria across incumbent strengths solves the equilibrium
again in each economy (Chapter~\ref{ch:ccc}).
\end{watchout}

\begin{checkyourself}
\label{q:auction-strict}
In the point-mass case with \(p<\ell<h\), find a condition on the strengthening \(\tilde F\le F\) under which
\(g_L\) does not change, \(g_H\) falls strictly, and \(\DT\) rises strictly. Explain the condition in words.
\end{checkyourself}

% ====================================================================
\section{Part (c): the weakest restriction}\label{sec:auction-weakest}

Part (b) assumes the class ranking. The appendix of the paper asks how much of the class ranking part (b) really
needs, and it answers with an exact condition.

\begin{paperresult}{Proposition 1(c) (the weakest restriction)}
Let \(G_H\) and \(G_L\) have finite means. The spread weakly increases under every first-order stochastic
strengthening of every continuous \(F\) if and only if \(D(y)\ge0\) for every \(y>p\) and \(D(p^-)\le0\). When
neither class has an atom at \(p\), this condition is exactly first-order dominance above the reserve together
with equal admissibility, \(G_H(p^-)=G_L(p^-)\).
\end{paperresult}

In plain words:
\begin{itemize}
\item \(D(y)\ge0\) for every \(y>p\) is \emph{dominance above the reserve}. Above the reserve, at every level, a
high-class value is at least as likely as a low-class value to lie above it. Below the reserve, this condition
puts no restriction on the order of the classes.
\item \(D(p^-)\le0\) means \(G_L(p^-)\le G_H(p^-)\). The low class fails the reserve at most as often as the high
class.
\item Without atoms at \(p\), the two conditions together say: dominance above the reserve, and equal
admissibility.
\end{itemize}
So the class ranking is stronger than part (b) needs. The class ranking asks for dominance at every level, and it
asks that neither class ever fail the reserve, \(G_H(p^-)=G_L(p^-)=0\). Without atoms at \(p\), part (c) asks
for dominance only above the reserve, and only that the two classes fail the reserve with the same probability.

\subsection*{The idea of the proof}

The proof works with \(\Gamma\), the class difference at one incumbent value (Section~\ref{sec:auction-slices}),
and with \(\DT(F)=\E_F[\Gamma(R)]\). A first-order strengthening weakly raises the average of a function that
never falls. Section~\ref{sec:toolkit-fosd} proves this for a differentiable function. The function \(\Gamma\) can jump at
\(p\), so Step~5 below does the algebra for \(\Gamma\) directly. So the spread weakly rises under every
strengthening when \(\Gamma\) never falls. Now suppose that \(\Gamma\) falls somewhere: \(\Gamma(R_1)>\Gamma(R_2)\)
at two values \(R_1<R_2\). A move of the incumbent from near \(R_1\) to near \(R_2\) is a strengthening, and it
lowers the spread. So the property holds exactly when \(\Gamma\) never falls, and the conditions on \(D\) say
exactly that.

\begin{derivation}[proof of Proposition~1(c)]
\begin{enumerate}[label=\textit{Step \arabic*.},leftmargin=*,itemsep=0.6em]
\item \emph{The kernel.} By Section~\ref{sec:auction-slices}, \(\DT(F)=\E_F[\Gamma(R)]\) with
\(\Gamma(R)=pD(p^-)\) for \(R<p\) and \(\Gamma(R)=\int_p^RD(y)\,dy\) for \(R\ge p\). The paper calls \(\Gamma\)
the kernel, and it notes that this kernel belongs to the cash second-price rule with reserve.
\item \emph{\(\Gamma\) is bounded.} Each \(G\) lies between zero and one, so
\begin{align*}
|D(y)|&=\Bigl|\bigl[1-G_H(y)\bigr]-\bigl[1-G_L(y)\bigr]\Bigr| \why{definition of \(D\)}\\
&\le\bigl[1-G_H(y)\bigr]+\bigl[1-G_L(y)\bigr] \why{\(|a-b|\le a+b\) for \(a,b\ge0\)}
\end{align*}
For \(R\ge p\), this gives
\begin{align*}
|\Gamma(R)|&\le\int_p^\infty|D(y)|\,dy \why{\(|\textstyle\int D|\le\int|D|\)}\\
&\le\int_0^\infty\bigl[1-G_H(y)\bigr]dy+\int_0^\infty\bigl[1-G_L(y)\bigr]dy
\why{the last bound}\\
&=\E_H[\theta]+\E_L[\theta]<\infty \why{tail formula; finite means}
\end{align*}
For \(R<p\), \(|\Gamma(R)|=p\,|D(p^-)|\le p\). So every expectation \(\E_F[\Gamma(R)]\) exists.
\item \emph{The shape of \(\Gamma\).} Below \(p\), \(\Gamma\) is constant at \(pD(p^-)\). At \(p\) it equals
\(\Gamma(p)=\int_p^pD\,dy=0\). Above \(p\), for \(p\le R_1<R_2\),
\begin{align*}
\Gamma(R_2)-\Gamma(R_1)&=\int_p^{R_2}D(y)\,dy-\int_p^{R_1}D(y)\,dy \why{definition of \(\Gamma\)}\\
&=\int_{R_1}^{R_2}D(y)\,dy \why{split the first integral at \(R_1\)}
\end{align*}
Since \(|D|\le1\), this increment is at most \(R_2-R_1\) in size. So \(\Gamma\) is continuous on \([p,\infty)\).
The paper also states that \(\Gamma\) has derivative \(D\) almost everywhere above \(p\). This guide works with
the increments \(\int_{R_1}^{R_2}D\,dy\) instead. They give the same conclusions and need no derivative.
\item \emph{When \(\Gamma\) never falls.} The paper calls this property \emph{nondecreasing}. The claim:
\(\Gamma\) is nondecreasing if and only if \(D(y)\ge0\) for every \(y>p\) and \(D(p^-)\le0\).
\begin{itemize}
\item Suppose the two conditions hold. On \([0,p)\), \(\Gamma\) is constant. At \(p\), it moves from
\(pD(p^-)\le0\) up to \(0\). On \([p,\infty)\), every increment \(\int_{R_1}^{R_2}D\,dy\) is nonnegative. So
\(\Gamma\) never falls.
\item Suppose \(\Gamma\) never falls. A value just below \(p\) gives \(pD(p^-)\le\Gamma(p)=0\), and \(p>0\), so
\(D(p^-)\le0\). Next suppose \(D(y_0)<0\) at some \(y_0>p\). Distribution functions are right-continuous, so
\(D\) is right-continuous. Then \(D(y)<\tfrac12D(y_0)<0\) on some interval \([y_0,y_0+\varepsilon)\), and
\(\Gamma(y_0+\varepsilon)-\Gamma(y_0)=\int_{y_0}^{y_0+\varepsilon}D\,dy<0\). This contradicts the assumption. So
\(D(y)\ge0\) for every \(y>p\).
\end{itemize}
\item \emph{The ``if'' direction.} Suppose \(D(y)\ge0\) for every \(y>p\) and \(D(p^-)\le0\). Take any continuous
\(F\) and any continuous \(\tilde F\le F\). The difference formula of Step~3 in the proof of Proposition~1 holds
for any classes:
\begin{align*}
\DT(\tilde F)-\DT(F)&=\int_p^\infty\bigl[F(y)-\tilde F(y)\bigr]D(y)\,dy\\
&\qquad-p\bigl[F(p)-\tilde F(p)\bigr]D(p^-) \why{difference formula}\\
&=\int_p^\infty\bigl[F(y)-\tilde F(y)\bigr]D(y)\,dy\\
&\qquad+p\bigl[F(p)-\tilde F(p)\bigr]\bigl(-D(p^-)\bigr) \why{move the minus sign}\\
&\ge0 \why{every factor \(\ge0\)}
\end{align*}
In the integral, both factors are nonnegative for \(y>p\). In the second term, \(p>0\),
\(F(p)-\tilde F(p)\ge0\), and \(-D(p^-)\ge0\). This is the algebra behind the general fact that a strengthening
weakly raises the average of a function that never falls.
\item \emph{The ``only if'' direction: two points.} Suppose the condition fails. By Step~4, \(\Gamma\) falls
somewhere: there are \(a<b\) with \(\Gamma(a)>\Gamma(b)\). The proof needs two such points \(0<R_1<R_2\), both
different from \(p\). They exist:
\begin{itemize}
\item \(\Gamma\) is constant on \([0,p)\), so \(b\ge p\). If \(a<p\), replace \(a\) by any point of \((0,p)\);
\(\Gamma\) keeps the same value there.
\item If \(a=p\), then \(\Gamma(p)>\Gamma(b)\) with \(b>p\). By continuity on \([p,\infty)\), a point \(a'\)
slightly above \(p\) still has \(\Gamma(a')>\Gamma(b)\). Use \(a'\).
\item If \(b=p\), then \(a<p\) and \(pD(p^-)>0=\Gamma(p)\). By continuity on \([p,\infty)\), a point \(b'\)
slightly above \(p\) still has \(\Gamma(b')<pD(p^-)\). Use \(b'\).
\item In every other case, \(a\) and \(b\) already work: take \(R_1=a\) and \(R_2=b\).
\end{itemize}
\item \emph{The ``only if'' direction: two incumbents.} Let \(F\) be uniform on \([R_1-\delta,R_1+\delta]\), and
let \(\tilde F\) be the same law shifted to \([R_2-\delta,R_2+\delta]\). Take \(\delta>0\) small, so that both
intervals lie in \((0,\infty)\) and neither contains \(p\). The shift gives
\(\tilde F(y)=F\bigl(y-(R_2-R_1)\bigr)\le F(y)\), because \(F\) never falls and \(R_2>R_1\). Both laws are
continuous. Then
\begin{align*}
\DT(\tilde F)-\DT(F)&=\E_{\tilde F}\bigl[\Gamma(R)\bigr]-\E_F\bigl[\Gamma(R)\bigr] \why{Step 1}\\
&\to\Gamma(R_2)-\Gamma(R_1)<0\quad\text{as }\delta\to0 \why{\(\Gamma\) continuous at \(R_1,R_2\)}
\end{align*}
Each expectation is the average of \(\Gamma\) over a short interval around \(R_1\) or \(R_2\). The average
lies between the smallest and the largest value of \(\Gamma\) on the interval. Since \(R_1,R_2\ne p\), \(\Gamma\)
is continuous near each point, so both bounds tend to the value at the point, and so does the average.
For a small enough \(\delta\), the difference is negative: this strengthening lowers the spread.
\item \emph{A support that starts at zero.} The model places \(F\) on \([0,\bar r]\). If the support must start at
zero, mix the same small uniform law \(U\) on \([0,\bar r]\) into both laws, with weight \(\varepsilon\in(0,1)\)
and \(\bar r\ge R_2+\delta\). Both mixtures are continuous, and \((1-\varepsilon)\tilde F+\varepsilon U\le
(1-\varepsilon)F+\varepsilon U\). The spread is an average of \(\Gamma\), so it is linear in the law:
\begin{align*}
&\DT\bigl((1-\varepsilon)\tilde F+\varepsilon U\bigr)-\DT\bigl((1-\varepsilon)F+\varepsilon U\bigr)\\
&\qquad=(1-\varepsilon)\E_{\tilde F}[\Gamma]+\varepsilon\E_U[\Gamma]-(1-\varepsilon)\E_F[\Gamma]-\varepsilon\E_U[\Gamma]
\why{linearity}\\
&\qquad=(1-\varepsilon)\bigl[\DT(\tilde F)-\DT(F)\bigr]<0 \why{Step 7}
\end{align*}
The mixture scales the difference by a positive factor and keeps its sign.
\item \emph{No atom at the reserve.} Right-continuity gives \(D(p)=\lim_{y\downarrow p}D(y)\ge0\) under
dominance above the reserve. Without atoms at \(p\), \(G_j(p^-)=G_j(p)\), so \(D(p^-)=D(p)\ge0\). Together with
\(D(p^-)\le0\), this forces \(D(p^-)=0\), which is \(G_H(p^-)=G_L(p^-)\): equal admissibility. In the other
direction, dominance above the reserve and equal admissibility give \(D(y)\ge0\) for \(y>p\) and \(D(p^-)=0\le0\).
\end{enumerate}
\end{derivation}

\begin{intuition}[part (c)]
The class difference in proceeds must never fall as the incumbent's value rises. Above the reserve, it rises by
\(D(y)\,dy\) at each level, so dominance above the reserve is needed. The one dangerous place is the reserve. Below
it, the difference is \(pD(p^-)\); just above it, the difference is zero. An incumbent that meets the reserve
protects the target against an inadmissible challenger: the incumbent then buys at \(p\). If the low class fails
the reserve more often, this protection helps the target more in class \(L\), and the class difference falls.
\end{intuition}

\begin{watchout}
The phrase ``when neither class has an atom at \(p\)'' matters in an edge case. With an atom at \(p\),
\(D(p^-)<0\) is possible while \(D(y)\ge0\) above the reserve. Then the reserve term \(pF(p)D(p^-)\) is negative,
a strengthening shrinks it toward zero, and the condition of part (c) holds without equal admissibility.
\end{watchout}

\subsection*{A reserve that excludes the low class}

Take point-mass classes with \(\ell<p<h\). The reserve lies between the two values. Then
\(G_L(p^-)=1\), since \(\ell<p\), and \(G_H(p^-)=0\). So \(D(p^-)=1>0\), and equal admissibility fails. Panel~(b)
of Figure~\ref{fig:auction-kernel} shows the result: \(\Gamma\) drops from \(p\) to zero at the reserve. An
incumbent that moves from just below the reserve to just above it lowers the spread. Question~\ref{q:auction-reserve}
works through the details.

The paper's main results keep the reserve below both classes, as in the class ranking. Chapter~\ref{ch:welfare}
treats reserves, among them a reserve that excludes the low class, where equal admissibility fails.

\subsection*{Scope of the class ranking}

The paper adds a scope remark. Its proofs use the class ranking only through two facts: \(D(y)\ge0\) for \(y>p\),
and \(D(p^-)=0\). These two facts give \(g_H\ge g_L\) and \(\DT\ge0\) (Section~\ref{sec:auction-spread}), the
comparison in part (b), and the class-gap identity of the next section. So the same results hold under dominance
above the reserve together with equal admissibility.

\begin{checkyourself}
\label{q:auction-reserve}
Take point-mass classes with \(\ell<p<h\). Find \(D(y)\) for \(y\ge p\), \(D(p^-)\) and \(\Gamma(R)\) for every
\(R\). Then let \(F\) be uniform on a short interval below \(p\), and let \(\tilde F\) be uniform on a short
interval \([R_2-\delta,R_2+\delta]\) with \(p<R_2-\delta\) and \(R_2+\delta<h\). Check that \(\tilde F\le F\). Compute
\(\DT(F)\) and \(\DT(\tilde F)\). For which \(R_2\) does the strengthening lower the spread? Explain in words.
\end{checkyourself}

% ====================================================================
\section{The class-gap identity}\label{sec:auction-gap}

The two formulas of Proposition~1 add up to a simple identity. It makes the story of
Section~\ref{sec:auction-story} exact.

\begin{derivation}[the sum of the two formulas]
Section~\ref{sec:auction-spread} found the profit gap \(g_H-g_L=\int_p^\infty D(y)F(y)\,dy\) from
\eqref{eq:auction-gross}; that step used no class ranking. Add the spread:
\begin{align*}
\DT+g_H-g_L&=pF(p)D(p^-)+\int_p^\infty\bigl[1-F(y)\bigr]D(y)\,dy\\
&\qquad+\int_p^\infty F(y)\,D(y)\,dy
\why{\eqref{eq:auction-spread}}\\
&=pF(p)D(p^-)+\int_p^\infty\bigl[1-F(y)+F(y)\bigr]D(y)\,dy \why{both integrals are finite}\\
&=pF(p)D(p^-)+\int_p^\infty D(y)\,dy \why{\(1-F+F=1\)}
\end{align*}
Next, the right side. For any value, \(\max\{\theta,p\}=p+(\theta-p)_+\): both sides equal \(p\) when
\(\theta\le p\) and \(\theta\) when \(\theta>p\). So
\begin{align*}
\E_j\bigl[\max\{\theta,p\}\bigr]&=p+\E_j\bigl[(\theta-p)_+\bigr] \why{linearity}\\
&=p+\int_p^\infty\Pr\nolimits_j(\theta>y)\,dy \why{tail formula}\\
&=p+\int_p^\infty\bigl[1-G_j(y)\bigr]\,dy \why{\(\Pr_j(\theta>y)=1-G_j(y)\)}
\end{align*}
The tail formula (Section~\ref{sec:toolkit-tail}) writes the mean of the nonnegative variable \((\theta-p)_+\) as
\(\int_0^\infty\Pr_j((\theta-p)_+>z)\,dz\). For \(z\ge0\), the event \((\theta-p)_+>z\) is the event
\(\theta>p+z\). Put \(y=p+z\); this is formula~\eqref{eq:toolkit-pospart} with \(c=p\). The difference of the classes:
\begin{align*}
\E_H\bigl[\max\{\theta,p\}\bigr]-\E_L\bigl[\max\{\theta,p\}\bigr]
&=\int_p^\infty\bigl[G_L(y)-G_H(y)\bigr]\,dy \why{\(p\) cancels}\\
&=\int_p^\infty D(y)\,dy \why{definition of \(D\)}
\end{align*}
Put the two results together:
\[
\DT(F)+g_H(F)-g_L(F)=pF(p)D(p^-)+\E_H\bigl[\max\{\theta,p\}\bigr]-\E_L\bigl[\max\{\theta,p\}\bigr].
\]
\end{derivation}

This is the paper's general form of the split. Under the class ranking, \(D(p^-)=0\), and the reserve term drops:
\begin{equation}\label{eq:auction-gap}
\DT(F)+g_H(F)-g_L(F)=\E_H\bigl[\max\{\theta,p\}\bigr]-\E_L\bigl[\max\{\theta,p\}\bigr].
\end{equation}
The reserve term also vanishes under equal admissibility, and when the incumbent always meets the reserve,
\(F(p)=0\). Under the class ranking every value meets the reserve, so \(\max\{\theta,p\}=\theta\) with
probability one, and the right side is the difference of the class means.

A second route to \eqref{eq:auction-gap} uses the split \eqref{eq:auction-split}. Under the class ranking,
\(\Pr_j(\theta\ge p)=1\), so \(t(\theta)+g(\theta)=\theta\) with probability one in each class. Take the class
expectation: \(t_j+g_j=\E_j[\theta]\). Subtract the low class from the high class:
\(\DT+g_H-g_L=\E_H[\theta]-\E_L[\theta]\).

\subsection*{What the identity says}

The right side of \eqref{eq:auction-gap} does not depend on \(F\). The paper says that the identity holds for
every incumbent distribution. In the model, every incumbent distribution is continuous, so the derivation above
covers each of them. (A remark of this guide: the second route does not use continuity, because the split holds
at every \(R\), ties included.) So a strengthening leaves the right side unchanged. Compare a strengthening
\(\tilde F\le F\):
\begin{align*}
\DT(\tilde F)-\DT(F)&=-\Bigl\{\bigl[g_H(\tilde F)-g_L(\tilde F)\bigr]-\bigl[g_H(F)-g_L(F)\bigr]\Bigr\}
\why{\eqref{eq:auction-gap} at \(\tilde F\) and at \(F\)}
\end{align*}
So the spread rises by exactly as much as the profit gap between the classes falls. A stronger incumbent leaves
the value gap between the classes unchanged. It moves the gap from the challenger's profit to the target's
proceeds.

In the slice picture, each slice \(D(y)\,dy\) goes either to the target or to the challenger. The incumbent's
distribution decides the shares, and the total stays fixed. The two extremes of Section~\ref{sec:auction-story}
are the two ends:
\begin{itemize}
\item An incumbent that never meets the reserve: \(\DT=0\), and \(g_H-g_L\) equals the whole gap.
\item An incumbent above every challenger value: \(g_H=g_L=0\), and \(\DT\) equals the whole gap.
\end{itemize}

\begin{watchout}
The identity concerns the profit gap \(g_H-g_L\). Each class profit \(g_j\) also falls under a strengthening, by
part (a), but the identity does not give the size of that fall. A strengthening can lower both \(g_H\) and \(g_L\)
by the same amount and leave \(\DT\) unchanged; the third case in the watch-out box of
Section~\ref{sec:auction-prop1} is an example.
\end{watchout}

% ====================================================================
\section{What the rest of the paper takes from this chapter}\label{sec:auction-uses}

The auction stage gives the rest of the paper the facts in the summary box below, and two more. Entry never
lowers proceeds, \(t_j\ge t_0\). Under the class ranking, \(g_H\ge g_L\) and \(\DT\ge0\). None of these facts
depends on a belief or a price.

The paper names Proposition~1 after the two returns to information. Gross profit \(g_j\), and through it \(B(\mu)\),
is what preparation earns the challenger once it learns its value. The spread \(\DT\) measures how much the
challenger's class matters to a holder of target shares. A strengthening weakly lowers the first return and, under
the class ranking, weakly raises the second. The later chapters turn these two returns into the objects that drive
behavior. This chapter does not prove those steps.
\begin{itemize}
\item Chapter~\ref{ch:beliefs} turns gross profit into the preparation threshold \(\tau\), the belief at which
\(B(\mu)=c_H\). When \(g_H>g_L\), \(\tau=(c_H-g_L)/(g_H-g_L)\). The paper notes that a stronger incumbent never
lowers the threshold while \(g_L\le c_H\).
\item Chapter~\ref{ch:trading} turns the spread into the return of the informed investor's trade. The gross return
of a correctly signed order is entry times the spread times the residual uncertainty about the class.
\end{itemize}
When the spread is positive and \(g_H>g_L\), the paper shows that orders, posteriors and entry depend on the
auction only through \(\DT\) and \(\tau\). Chapter~\ref{ch:ccc} then weighs the lower profit against the more
informative price.

\begin{chaptersummary}
\begin{itemize}
\item Each prepared bidder whose value meets the reserve bids its value. So the winner pays the larger of the
reserve and the value of the other bidder, when that bidder is admissible.
\item For a challenger value \(\theta\ge p\), the target receives \(\max\{p,\min(R,\theta)\}\) and the challenger
keeps \((\theta-\max\{p,R\})_+\); the two add up to \(\theta\), so \(t(\theta)+g(\theta)=\theta\).
\item Written as sums over levels, \(g_j=\int_p^\infty[1-G_j(y)]F(y)\,dy\) and
\(\DT=pF(p)D(p^-)+\int_p^\infty[1-F(y)]D(y)\,dy\).
\item A strengthening lowers \(F\). So it weakly lowers \(g_H\) and \(g_L\) for every pair of classes, and under
the class ranking it weakly raises \(\DT\). A strict change needs the stated integral to be positive.
\item When the classes are ranked above the reserve, the reserve term is the only channel through which a
strengthening can narrow the spread. Part (c) shows that, without atoms at \(p\), dominance above the reserve and
equal admissibility are exactly what part (b) needs.
\item Under the class ranking, \(\DT+g_H-g_L\) equals the class value gap, which does not depend on the incumbent:
a stronger incumbent moves the gap from the challenger's profit to the target's proceeds.
\end{itemize}
\end{chaptersummary}

% ====================================================================
\answersection

\answer{\ref{q:auction-overbid}}{%
Both bids are admissible, since \(b>\theta\ge p\). Go through the three ranges of \(Y\).
\begin{itemize}
\item \(Y<\theta\). Both bids are above \(Y\), so both win. Both pay \(\max\{p,Y\}\), which does not depend on the
bid. Both earn \(\theta-\max\{p,Y\}\). The payoffs are equal.
\item \(Y>b\). Both bids are below \(Y\), so both lose and earn zero. The payoffs are equal.
\item \(\theta\le Y\le b\). The bid \(\theta\) loses, or ties at \(Y=\theta\) and then pays \(\theta\). Either way
it earns zero. The bid \(b\) wins, or ties at \(Y=b\). If it wins, it pays \(\max\{p,Y\}=Y\), since \(Y\ge\theta\ge
p\), and it earns \(\theta-Y\le0\). If it loses, it earns zero. So the bid \(b\) earns at most zero.
\end{itemize}
The bid \(b\) never earns more than the bid \(\theta\). It is strictly worse when \(\theta<Y<b\): it then wins and
pays \(Y>\theta\), a loss of \(Y-\theta\). It is also strictly worse at \(Y=b\) when the tie goes to the bid \(b\),
for the same reason.}

\answer{\ref{q:auction-jump}}{%
At \(\theta=p\), the target receives \(\max\{p,\min(R,p)\}\) at every \(R\). Since \(\min(R,p)\le p\), this is
\(p\). The challenger keeps \((p-\max\{p,R\})_+=0\), since \(\max\{p,R\}\ge p\). So
\begin{align*}
t(p)&=\E_R[p]=p \why{the proceeds are \(p\) at every \(R\)}\\
g(p)&=\E_R[0]=0 \why{the profit is zero at every \(R\)}
\end{align*}
The split holds: \(t(p)+g(p)=p\). For \(\theta<p\), \(t(\theta)=t_0=p[1-F(p)]\). The difference is
\begin{align*}
t(p)-t_0&=p-p\bigl[1-F(p)\bigr] \why{the two values}\\
&=pF(p) \why{expand}
\end{align*}
In words: when \(R\ge p\), the target receives \(p\) in both cases. A challenger at \(p\) loses to the incumbent,
which pays \(\max\{p,\theta\}=p\), and a challenger below \(p\) leaves the incumbent to pay \(p\). When \(R<p\),
which has probability \(F(p)\), the two cases differ. A challenger at \(p\) is admissible and buys at \(p\). A
challenger below \(p\) cannot bid, and the target stays unsold. This difference is the jump \(pF(p)\) in
Section~\ref{sec:auction-spread}.}

\answer{\ref{q:auction-pointgross}}{%
With point masses, \(G_H(y)=\1\{y\ge h\}\) and \(G_L(y)=\1\{y\ge\ell\}\). So \(1-G_H(y)=\1\{y<h\}\) and
\(1-G_L(y)=\1\{y<\ell\}\). Then
\begin{align*}
g_H&=\int_p^\infty\1\{y<h\}\,F(y)\,dy=\int_p^hF(y)\,dy \why{\eqref{eq:auction-gross}}\\
g_L&=\int_p^\infty\1\{y<\ell\}\,F(y)\,dy=\int_p^\ell F(y)\,dy \why{\eqref{eq:auction-gross}}\\
g_H-g_L&=\int_\ell^hF(y)\,dy \why{\(p<\ell<h\); split \([p,h]\) at \(\ell\)}
\end{align*}
In words, the integral counts, at each level \(y\) between \(\ell\) and \(h\), the probability that the incumbent is below
\(y\). A high-class challenger keeps the slice at such a level when it beats the incumbent there. A low-class
challenger never reaches these levels. Since \(F\ge0\), the integral is nonnegative, so \(g_H\ge g_L\). It is zero
only when \(F(y)=0\) on \([\ell,h]\): then the incumbent is always at or above \(h\), and both classes always lose.}

\answer{\ref{q:auction-slices}}{%
Fix \(R\) with \(\ell\le R<h\). Both class values meet the reserve, since \(p<\ell<h\).
\begin{itemize}
\item Class \(H\): \(h>R\), so the challenger wins and pays \(\max\{p,R\}=R\). The target receives \(R\), and the
challenger keeps \(h-R\).
\item Class \(L\): \(\ell\le R\), so the incumbent wins, or ties at \(R=\ell\), and the price is \(\ell\). The
target receives \(\ell\), and the challenger keeps \((\ell-R)_+=0\).
\end{itemize}
The class difference in proceeds is \(R-\ell\). The formula agrees:
\(\Gamma(R)=\int_p^R\1\{\ell\le y<h\}\,dy=R-\ell\), since \(\ell\le R<h\). The class difference in the challenger's
profit is \((h-R)-0=h-R\). Add the two:
\begin{align*}
\Gamma(R)+(h-R)&=(R-\ell)+(h-R) \why{the two differences}\\
&=h-\ell \why{\(R\) cancels}
\end{align*}
This is the class value gap \(\E_H[\max\{\theta,p\}]-\E_L[\max\{\theta,p\}]=h-\ell\). The slices between \(\ell\)
and \(R\) reach the target: the high-class challenger pays for them, and the low class does not have them. The
slices between \(R\) and \(h\) stay with the high-class challenger as profit. A higher \(R\) moves the boundary
line toward \(h\) and sends more slices to the target.}

\answer{\ref{q:auction-strict}}{%
With point masses, \(1-G_L(y)=\1\{y<\ell\}\), \(1-G_H(y)=\1\{y<h\}\) and \(D(y)=\1\{\ell\le y<h\}\). Proposition~1
and its proof give the three changes:
\begin{align*}
g_L(F)-g_L(\tilde F)&=\int_p^\ell\bigl[F(y)-\tilde F(y)\bigr]dy \why{Step 1, class \(L\)}\\
g_H(F)-g_H(\tilde F)&=\int_p^h\bigl[F(y)-\tilde F(y)\bigr]dy \why{Step 1, class \(H\)}\\
\DT(\tilde F)-\DT(F)&=\int_\ell^h\bigl[F(y)-\tilde F(y)\bigr]dy \why{Step 4}
\end{align*}
The function \(F-\tilde F\) is continuous and nonnegative. Such a function has a zero integral over an interval
only when it is zero on the whole interval. So:
\begin{itemize}
\item \(g_L\) does not change exactly when \(\tilde F(y)=F(y)\) for every \(y\in[p,\ell]\).
\item Given that, \(g_H\) falls strictly exactly when \(\int_\ell^h[F-\tilde F]\,dy>0\), that is, when
\(\tilde F(y)<F(y)\) at some \(y\in(\ell,h)\).
\item The same condition makes \(\DT\) rise strictly.
\end{itemize}
The condition: \(\tilde F=F\) on \([p,\ell]\), and \(\tilde F<F\) somewhere in \((\ell,h)\). In words, the
strengthening keeps the chance that the incumbent is below each level up to \(\ell\). It moves some incumbent
values that lie above \(\ell\) to higher values. The low class loses to every incumbent above \(\ell\)
anyway, so nothing changes for it. The high class meets a higher price or loses more often, so its profit falls.
In class \(L\) the target still receives \(\ell\) from those incumbents, while in class \(H\) it receives more. So
the spread rises by exactly the fall in \(g_H\), as the class-gap identity requires.}

\answer{\ref{q:auction-reserve}}{%
With \(\ell<p<h\): \(G_L(y)=\1\{y\ge\ell\}\) and \(G_H(y)=\1\{y\ge h\}\).
\begin{itemize}
\item \(G_L(p^-)=1\), since \(\ell<p\), and \(G_H(p^-)=0\), since \(h>p\). So \(D(p^-)=1\).
\item For \(y\ge p\): \(G_L(y)=1\), so \(D(y)=1-\1\{y\ge h\}=\1\{y<h\}\).
\item For \(R<p\): \(\Gamma(R)=pD(p^-)=p\). For \(R\ge p\):
\(\Gamma(R)=\int_p^R\1\{y<h\}\,dy=\min\{R,h\}-p\).
\end{itemize}
First, \(\tilde F\le F\): for \(y\ge p\), \(F(y)=1\ge\tilde F(y)\), and for \(y<p\), \(\tilde F(y)=0\le F(y)\).
The incumbent \(F\) lies below \(p\), so \(\Gamma(R)=p\) at every value it takes. The incumbent \(\tilde F\) lies in
\((p,h)\), so \(\Gamma(R)=R-p\) at every value it takes. Then
\begin{align*}
\DT(F)&=\E_F[\Gamma(R)]=p \why{\(\Gamma=p\) below \(p\)}\\
\DT(\tilde F)&=\E_{\tilde F}[R-p]=R_2-p \why{the mean of \(\tilde F\) is \(R_2\)}\\
\DT(\tilde F)-\DT(F)&=R_2-2p \why{subtract}
\end{align*}
So the strengthening lowers the spread when \(R_2<2p\), together with the stated \(p<R_2-\delta\) and
\(R_2+\delta<h\). Neither law starts at zero; Step~8 of the proof of part (c) mixes in a small uniform law on
\([0,\bar r]\), and this keeps the sign.

In words: when the incumbent fails the reserve, the high-class challenger buys at \(p\), and the low-class
challenger cannot bid, so the target stays unsold. The class difference is \(p\). When the incumbent meets the
reserve, it buys at \(p\) whenever the challenger is in the low class, and the high class pays about \(R_2\). The
class difference is about \(R_2-p\), which is small when \(R_2\) is near \(p\). The stronger incumbent protects the
target against the excluded low class, and that protection narrows the class difference. Equal admissibility
fails here, and part (c) says some strengthening must lower the spread; this is one.}
